\chapter{虚拟内存}

本章对应原书第~9~章 Virtual Memory。课堂尚未讲授，以下仅保留目录。

\section{物理和虚拟寻址}

\pending

\section{地址空间}

\pending

\section{虚拟内存作为缓存的工具}

\pending

\subsection{DRAM缓存的组织}

\pending

\subsection{页表}

\pending

\subsection{页命中}

\pending

\subsection{缺页}

\pending

\subsection{分配页面}

\pending

\subsection{再论局部性}

\pending

\section{虚拟内存作为内存管理的工具}

\pending

\section{虚拟内存作为内存保护的工具}

\pending

\section{地址翻译}

\pending

\subsection{结合缓存和虚拟内存}

\pending

\subsection{利用TLB加速地址翻译}

\pending

\subsection{多级页表}

\pending

\subsection{端到端的地址翻译}

\pending

\section{案例研究：Intel Core i7/Linux内存系统}

\pending

\subsection{Core i7地址翻译}

\pending

\subsection{Linux虚拟内存系统}

\pending

\section{内存映射}

\pending

\subsection{再看共享对象}

\pending

\subsection{再看fork函数}

\pending

\subsection{再看execve函数}

\pending

\subsection{使用mmap函数的用户级内存映射}

\pending

\section{动态内存分配}

\pending

\subsection{malloc和free函数}

\pending

\subsection{为什么要使用动态内存分配}

\pending

\subsection{分配器的要求和目标}

\pending

\subsection{碎片}

\pending

\subsection{实现问题}

\pending

\subsection{隐式空闲链表}

\pending

\subsection{放置已分配块}

\pending

\subsection{分割空闲块}

\pending

\subsection{获取额外的堆内存}

\pending

\subsection{合并空闲块}

\pending

\subsection{带边界标记的合并}

\pending

\subsection{实现一个简单的分配器}

\pending

\subsection{显式空闲链表}

\pending

\subsection{分离的空闲链表}

\pending

\section{垃圾收集}

\pending

\subsection{垃圾收集器的基本知识}

\pending

\subsection{Mark\&Sweep垃圾收集器}

\pending

\subsection{C程序的保守Mark\&Sweep}

\pending

\section{C程序中常见的与内存有关的错误}

\pending

\subsection{间接引用坏指针}

\pending

\subsection{读未初始化的内存}

\pending

\subsection{允许栈缓冲区溢出}

\pending

\subsection{假设指针和它们指向的对象是相同大小的}

\pending

\subsection{造成错位错误}

\pending

\subsection{引用指针，而不是它所指向的对象}

\pending

\subsection{误解指针运算}

\pending

\subsection{引用不存在的变量}

\pending

\subsection{引用空闲堆块中的数据}

\pending

\subsection{引起内存泄漏}

\pending

\section{小结}

\pending
